% Generated draft: requires mathematical and source review.
\documentclass[11pt]{article}
\usepackage[margin=25mm]{geometry}
\usepackage{fontspec}
\setmainfont{Latin Modern Roman}
\usepackage{amsmath,amssymb,mathtools,unicode-math}
\setmathfont{Latin Modern Math}
\usepackage{enumitem,needspace,xurl,hyperref}
\setlength{\parindent}{0pt}
\setlength{\parskip}{5pt}
\setlength{\emergencystretch}{3em}
\newcommand{\N}{\mathbb{N}}
\newcommand{\Z}{\mathbb{Z}}
\newcommand{\Q}{\mathbb{Q}}
\newcommand{\R}{\mathbb{R}}
\newcommand{\C}{\mathbb{C}}
\newcommand{\F}{\mathbb{F}}
\newcommand{\A}{\mathbb{A}}
\newcommand{\PP}{\mathbb{P}}
\newcommand{\E}{\mathbb{E}}
\newcommand{\eps}{\varepsilon}
\newcommand{\dd}{\,\mathrm{d}}
\begin{document}
\section*{UnsolvedMath ID: 6700025}
\textbf{Source problem code:} catalogue entry
\subsection{Definitions and mathematical statement}
Let \(U(N)=\{A\in M_N(\C):A^*A=I\}\). Once metrics \(d_S\) on \(S^{n-1}\) and \(d_U\) on \(U(N)\) are fixed, define
\[
\operatorname{Lip}(\Phi)=\sup_{x\ne y}\frac{d_U(\Phi(x),\Phi(y))}{d_S(x,y)},\qquad
\mathcal L_{n,N}=\{\Phi:S^{n-1}\to U(N):\operatorname{Lip}(\Phi)<1/2\}.
\]
The threshold \(1/2\) is the likely fraction reading of the extracted ``1 2'' and requires confirmation. Gromov suggests that these maps may be deformable to constant maps \emph{continuously in \(\Phi\)}, for all \(n\). A parameter-dependent null-homotopy formulation would ask for a continuous map
\[
 H:\mathcal L_{n,N}\times[0,1]\longrightarrow C(S^{n-1},U(N))
\]
with \(H(\Phi,0)=\Phi\) and \(H(\Phi,1)\) constant on \(S^{n-1}\) for every \(\Phi\). The topology on the map spaces must be specified; the source does not state it. It also does not say whether \(H(\Phi,t)\) must remain in \(\mathcal L_{n,N}\), satisfy some other Lipschitz bound, or fix constant maps throughout.

This asks for a continuously chosen deformation, not merely a separate null-homotopy for each map. Nor does deforming to the \emph{space} of constant maps assert contractibility to one fixed constant map. The preceding argument and its ``last step'', which motivates the suggestion, are not supplied. The metrics, the range of \(n\), and whether the assertion is uniform in \(N\) are likewise unstated.
\subsection{Short English statement}
Can sufficiently small Lipschitz maps from spheres to unitary groups be deformed to constant maps by a choice that varies continuously with the original map, in every dimension? Metric normalizations and the required topology and Lipschitz control of the deformation are missing.
\subsection{Sources}
\begin{itemize}
\item Local UnsolvedMath registry, permanent ID 6700025, source code catalogue entry; Misha Gromov, 2017; update 2026-08-06.
\item Gromov, \emph{101 Questions, Problems and Conjectures around Scalar Curvature} (2017), printed tag [?24], occurrence 25, PDF page 36, passage (i): \url{https://www.ihes.fr/~gromov/wp-content/uploads/2018/08/101-problemsOct1-2017.pdf}. Registry access date 2026-07-29; normalized PDF extraction; incomplete public author-hosted PDF; status and reuse terms require review.
\item The supplied summary associates the question with a K-area/index argument and Lipschitz-topology contractibility but gives no additional reference or precise formulation. Its description as involving high Lipschitz constants does not match the extracted small-threshold inequality.
\end{itemize}
\end{document}
